\noindent \textit{\textcolor{red!60!black}{
One small point on last sentence page 12/top page 13. I don't think that the Liberty bond prices reflect a flight to safety. Indeed, Liberty bonds were very much part of debate since they were to be paid in gold and in the end were not due to the US leaving the gold standard (which is why some argue that this constitutes sovereign default). If anything, the high prices of Liberty bonds indicated that observations thought it likely that the government could lose the case. Either the authors provide a better explanation for the Liberty bond behavior, which is currently stated without the appropriate context or they delete the sentence. Also, the citation to the NYTimes article is missing.
}}

\bigskip

\noindent Thank you for this comment. We have revised the passage to provide the appropriate context for the Liberty Bond price increase. Liberty Bonds were themselves gold-denominated government obligations, promising payment in gold at par. A government loss in the gold clause cases would have obligated the Treasury to honor that promise at the then-prevailing price of \$35 per ounce---69\% above the pre-abrogation level---increasing the real value of the bonds to holders. Their rising prices in January 1935 therefore reflected investor expectations that the government could lose the case. The historical background section now reads:

\begin{quote}
``During the three days of oral arguments (January 8--10), the market fell 1.3\% and our gold-exposure-weighted portfolio declined 1.4\%. These modest declines understate the anxiety the arguments triggered, in part because the news spread slowly: many newspapers reported on the government's poor showing only on January 12 and 13 \citep{Edwards2018}. As the news diffused, the market fell a further 3.0\% between January 11 and January 17, while the gold-exposure-weighted portfolio declined 4.6\%---cumulative declines of 4.3\% and 6.0\%, respectively, since the start of the arguments. By January 13, the \textit{New York Times} reported that, amid `the speculative fever for gold,' Liberty Bond prices had reached their highest level since 1917 \citep{NYT1935}. Since Liberty Bonds were themselves gold-denominated government obligations, their rising prices reflected investor expectations that the government could lose the case and be forced to honor full gold payments.''
\end{quote}

We have also added the missing citation: \textit{New York Times}, ``Gold Issue Hits Highest since 1917,'' January 13, 1935, p.\ 2.
